\section{Introduction}

This document is a curated portfolio of algorithmic solutions.
Each selected problem includes:

\begin{itemize}[nosep]
  \item clear problem restatements,
  \item key invariants / observations,
  \item pseudocode,
  \item correctness sketches,
  \item complexity analysis,
  \item C++ (primary) and Python (secondary) implementations.
\end{itemize}

The focus is on clarity, correctness, and reasoning.


\subsection{How to Use This Book}

Problems include an optional \textbf{Author's Thinking} block, the authors plain language thinking and problem framing, which is excluded from exemplar builds.
Use this book to understand how to formally reason with and structure algorithm and data structures problems.


\subsection{Motivation}

This portfolio was created as a structured, self-directed study of fundamental
algorithms and data structures. While my formal coursework emphasizes applied
computer science and systems-level engineering, it does not include a dedicated
course focused on algorithmic problem-solving in the style commonly expected in
technical interviews and quantitative software roles.

To address this, I undertook a systematic review of core algorithmic techniques,
organizing problems by data structure and algorithmic paradigm, and documenting
each solution with an explicit algorithm description, correctness argument, and
complexity analysis. The goal is not only to arrive at correct implementations,
but to practice communicating algorithmic reasoning clearly and rigorously.

This document serves both as a personal reference and as evidence of deliberate
practice in algorithmic reasoning beyond formal coursework.

Many of the problems included here are sourced from commonly used interview
problem sets, and were selected to emphasize invariant-based reasoning, edge
case analysis, and asymptotic performance trade-offs.

Informal reasoning notes are included selectively to document thought process
and learning progression; these sections are omitted in derived exemplar or
presentation-focused versions of this material.


\subsection{Acknowledgements}

This work was informed by a number of open educational resources that emphasize
rigorous algorithmic reasoning and clear exposition. In particular, the structure
and approach to correctness arguments and asymptotic analysis were influenced by
\emph{MIT OpenCourseWare 6.006: Introduction to Algorithms}.

Additional reference was made to publicly available problem discussions,
editorials, and textbooks where appropriate. All implementations, explanations,
and written analyses in this document are my own, and any external influences are
acknowledged at a conceptual level rather than reproduced directly.
